% !TeX root = ../papers/nc-p3.tex
\section{四次以下に分かれる反対八次式と四面体}
\label{sec:quartic-octic-splitting}
\begin{theorem}\label{thm:quartic-octic-splitting}
定理\ref{thm:reducible-octic-factors}と同じ仮定で、八次式の既約成分がすべて四次以下なら、
一つの既約四次曲面と四面体の四面に分かれる場合だけが残る。四頂点は印で四次曲面の外にあり、
残る十六点は四次曲面の二重点で、各面に四点ずつ載り、辺を避ける。
特に既約な二つの四次曲面への分解は生じない。
\end{theorem}
\begin{proof}[低重複度の因子の核座標]
成分の既約な因子 $F$ とその曲面 $S$ を取り、$S'$ を点の吹上げでの真変換とする。
$S$ が印で三重にならなければ、$S'$ 上で $L=P$ となる。これを特異な $S$ に対しても示す。
$\pi:S'\to S$ の幾何学的繊維は、一点か、被約な一次または二次の平面曲線である。
実際印での $F$ の初期形式は被約な $d$ の初期形式の因子である。これらの繊維は連結で
$H^1_{\text{ét}}(-,\mathbb Z/2)=0$（二直線なら交点で貼り合わせる）。
固有基底変換とLerayにより $H^1_{\text{ét}}(S',\mathbb Z/2)=H^1_{\text{ét}}(S,\mathbb Z/2)=0$。
後の消滅は特異な超曲面にも成立する。実際複素数上に移し、$U=\mathbb P^3\setminus S$ は滑らかな
アフィン三次元多様体ゆえ実次元3以下のCW複体のホモトピー型を持つ（Andreotti--Frankel）。
従って $H^2_c(U,\mathbb Z/2)=H_4(U,\mathbb Z/2)=0$ と $H^1(\mathbb P^3,\mathbb Z/2)=0$ から従う。
有限個の係数の定義体を用い、比較定理と固有基底変換で元の標数零代数閉体にも戻る。
Kummer列により $L-P$ は自明である。また $S$ は超曲面ゆえ $S_2$ を満たし、有限点を除く部分の正則関数は延長する。
従って $\pi_*\mathcal O_{S'}=\mathcal O_S$ で、$L$ の切断は印で消える周囲の三次式から来る。
特に四次成分で $n=|X\cap S|$ なら $h^0(L)\le20-n$ である。

次数4の成分の核像は直線に入らない。入ればその固定二空間の基底に対応する $N$ の二列の
$2\times2$ 小行列式は全て $F$ の定数倍である。二列は一般点で独立なので、その張る平面の
Plücker座標が入力全体で一定になる。$N(x)$ の値は全対称行列空間を張り、固定した非零ベクトルの像は
四空間を張るので矛盾する。従って $h^0(L)\ge3$。
さらに成分が二つの異なる印で二重なら核像は平面にも入らない。例外平面上の六つの余因子は単元の消去後に
全二次式の基底となるため、その例外平面からの余因子射は閉埋込みである。接二次曲線上では出力のVeroneseの逆を
経由し、核座標の次数は各成分で $\mathcal O(1)$ と等しい（その二乗で比較する）。
従って被約な二次曲線は閉埋込みにより非共線な二次曲線となり、その張る平面は $y(p)^\perp$ である。
二点で出力印 $y(p)$ が異なるから、二つの平面は一致しない。この場合 $h^0(L)\ge4$。
\end{proof}
\begin{proof}[四次曲面の三重点からの射影]
既約四次曲面が印 $p$ で三重なら、その上の特異な印は $p$ を含め高々七個である。
$p$ を頂点として方程式を $w f_3(u)+f_4(u)=0$ と書く。$f_3$ は接三次で非零、
既約性から $f_3,f_4$ に共通因子はない。$p$ 以外の特異点の方向では両式が零で、
$w\,df_3+df_4=0$。両微分が零でなければその方向の特異点は高々一点で、局所交点数は2以上である。
両微分が零なら交点数は4以上で、その直線の $p$ 以外の印は高々二個である。
従って残りの特異な印の数は $3\cdot4/2=6$ 以下となる。

既約四次の三重の印は高々二個である。まず任意の四次式で三つの三重の印が共線なら、
その直線は三重直線となる。全二階偏導二次式がこの直線を三点で消えるからである。
このとき三重点での接三次式は横断二変数の三次式であり、三共点直線か重複を持つ。
いずれも対称表示の非自明な二捩れと両立しない。被約三共点直線の
Picard群の次数零部分は加法群である。実際正規化の三枝の定数貼合せは各成分の定数で消え、
残る単元商は対数により単冪部分である。四次式が二つの三重点を持てばその間の直線は二重である
（一次偏導三次式が両端で二重）。従って非共線な三つの三重点は、三辺を各二重に含む平面切断を作り、
その平面を因子に強制する。これは既約四次では不可能である。以上より三重の印を持つ既約四次の重複度の和は
$20+7+2=29$ 以下である。

既約四次が二つの場合、二十点はどれも少なくとも片方で特異である。双方に三重の印があれば
特異な印の合計は高々14。三重の印のない側があれば $3\le h^0(L)\le20-n$ なので
他方に三重の印が少なくとも三つ必要となる。いずれも矛盾する。
\end{proof}
\begin{proof}[残余四次式の四つの平面]
定理\ref{thm:reducible-octic-factors}により既約四次 $F$ が一つある。
残余四次 $G$ の全因子が三次以下なら、それらの重複度の和は高々 $22+7=29$。
$F$ の和は31以上なので $F$ に三重の印はなく、二重の印が二つ以上存在する。
従って上の核像の議論で $4\le h^0(L)\le20-n$ となり、$F$ の外で $G$ は四点以上を三重に通る。
そのうちのどの三点も共線でなく、それらの平面は $G$ の因子となる。
四点が同一平面上にあればこの因子を割った三次式はその四点を二重に通り、
同平面で六割線を含むためもう一度同平面で割れる。平方自由性に反する。
四非共面点なら四つの三点平面で $G$ が尽くされる。その三重点は四頂点に限られる。

全二十点は $G$ に載る。四面の印の個数の和は $4\cdot7=28$ 以下で、
四頂点を各三回、他の十六点を一回数えるだけで既に28となる。
よって各面に七個、四頂点以外は辺上になく、各面で頂点を除く四点では $F$ が二重である。
\end{proof}
